\subsection{ベクトルと基底}

大きさだけをもつ量を\textbf{スカラー}、大きさと向きをもつ量を
\textbf{ベクトル}という。ベクトルは太字の $\vect{A}$ で表し、
その大きさを $|\vect{A}|$ と書く。

\subsubsection{基底による表現}

$n$ 次元空間の基底を $\{\vect{a}_1,\vect{a}_2,\ldots,\vect{a}_n\}$
とすると、任意のベクトル $\vect{A}$ は
\begin{equation}
  \vect{A}=x_1\vect{a}_1+x_2\vect{a}_2+\cdots+x_n\vect{a}_n
  =\begin{pmatrix}\vect{a}_1&\vect{a}_2&\cdots&\vect{a}_n\end{pmatrix}
   \begin{pmatrix}x_1\\x_2\\\vdots\\x_n\end{pmatrix}
\end{equation}
と表される。この係数の組を、基底に関する $\vect{A}$ の\textbf{成分}という。

三次元の正規直交基底 $\{\unit{x},\unit{y},\unit{z}\}$ は
\begin{equation}
  \unit{i}\cdot\unit{j}=\delta_{ij},\qquad
  \delta_{ij}=\begin{cases}1&(i=j),\\0&(i\ne j)\end{cases}
\end{equation}
を満たす。ここで $i,j\in\{x,y,z\}$ である。したがって
\begin{equation}
  \vect{A}=A_x\unit{x}+A_y\unit{y}+A_z\unit{z}
  \quad\Longleftrightarrow\quad
  \vect{A}=\begin{pmatrix}A_x\\A_y\\A_z\end{pmatrix}.
\end{equation}

\begin{figure}[htbp]
  \centering
  \begin{tikzpicture}[>=Stealth,scale=1.05]
    \draw[->] (0,0)--(4,0) node[right] {$x$};
    \draw[->] (0,0)--(0,3) node[above] {$y$};
    \draw[->,very thick] (0,0)--(3.1,2.1) node[above right] {$\vect{A}$};
    \draw[dashed] (3.1,0) node[below] {$A_x$}--(3.1,2.1);
    \draw[dashed] (0,2.1) node[left] {$A_y$}--(3.1,2.1);
    \draw[->,thick] (0,0)--(0.8,0) node[below] {$\unit{x}$};
    \draw[->,thick] (0,0)--(0,0.8) node[left] {$\unit{y}$};
  \end{tikzpicture}
  \caption{正規直交基底によるベクトルの成分表示}
\end{figure}

\subsubsection{スカラー場とベクトル場}

空間の各点 $P$ に一つの数 $\phi(P)$ を対応させたものを
\textbf{スカラー場}という。温度や密度がその例である。
スカラー場の変化率は勾配
\begin{equation}
  \nabla\phi=
  \frac{\partial\phi}{\partial x}\unit{x}
  +\frac{\partial\phi}{\partial y}\unit{y}
  +\frac{\partial\phi}{\partial z}\unit{z}
\end{equation}
で表される。勾配は、値が最も速く増加する向きを向き、その大きさは
その方向の変化率を表す。

\begin{figure}[htbp]
  \centering
  \begin{tikzpicture}[>=Stealth,scale=.95]
    \draw[rounded corners,fill=blue!4] (0,0) rectangle (5,3.1);
    \foreach \r in {.35,.65,1.0}
      \draw[blue!55] (2.4,1.6) ellipse ({\r*1.45} and {\r});
    \fill (1.45,1.15) circle (1.5pt) node[below left] {$P$};
    \node at (2.4,1.6) {$\phi$ が大きい};
    \draw[->,very thick,red!70!black] (1.45,1.15)--(2.65,1.75)
      node[above] {$\nabla\phi$};
    \node[align=left] at (7.3,2.2) {各点に数 $\phi(P)$\\$\Rightarrow$ スカラー場};
    \node[align=left] at (7.3,.8) {最大増加方向\\$\Rightarrow$ 勾配ベクトル};
  \end{tikzpicture}
  \caption{スカラー場とその勾配}
\end{figure}

\subsection{ベクトルの演算}

ベクトルの和は平行四辺形の対角線として幾何学的に表される。
成分表示では、対応する成分どうしを加えればよい。
\begin{equation}
  \vect{A}+\vect{B}
  =\begin{pmatrix}A_x+B_x\\A_y+B_y\\A_z+B_z\end{pmatrix}.
\end{equation}

\begin{figure}[htbp]
  \centering
  \begin{tikzpicture}[>=Stealth]
    \coordinate (O) at (0,0); \coordinate (A) at (1.4,2.2);
    \coordinate (B) at (3.4,.6); \coordinate (C) at ($(A)+(B)$);
    \draw[->,very thick] (O)--(A) node[midway,left] {$\vect{A}$};
    \draw[->,very thick] (O)--(B) node[midway,below] {$\vect{B}$};
    \draw[dashed] (A)--(C)--(B);
    \draw[->,very thick] (O)--(C) node[midway,above left] {$\vect{A}+\vect{B}$};
  \end{tikzpicture}
  \caption{ベクトルの加法（平行四辺形則）}
\end{figure}

加法とスカラー倍には、次の法則が成り立つ。
\begin{align}
  \vect{A}+\vect{B}&=\vect{B}+\vect{A}, &
  \vect{A}+(\vect{B}+\vect{C})&=(\vect{A}+\vect{B})+\vect{C},\\
  (ab)\vect{A}&=a(b\vect{A}), &
  (a+b)\vect{A}&=a\vect{A}+b\vect{A},\\
  a(\vect{A}+\vect{B})&=a\vect{A}+a\vect{B}.&&
\end{align}

\begin{example}[内分点の位置ベクトル]
位置ベクトル $\vect{r}_1,\vect{r}_2$ の点を結ぶ線分を
$m_2:m_1$ に内分する点の位置ベクトルは
\begin{equation}
  \vect{R}=\vect{r}_1+\frac{m_2}{m_1+m_2}(\vect{r}_2-\vect{r}_1)
  =\frac{m_1\vect{r}_1+m_2\vect{r}_2}{m_1+m_2}
\end{equation}
である。
\begin{center}
  \begin{tikzpicture}[>=Stealth,scale=.9]
    \coordinate (O) at (0,0); \coordinate (P1) at (2,2.4);
    \coordinate (P2) at (6,1.1); \coordinate (P) at (3.7,1.85);
    \draw[->] (O)--(P1) node[midway,left] {$\vect{r}_1$};
    \draw[->] (O)--(P2) node[midway,below] {$\vect{r}_2$};
    \draw[->,very thick] (O)--(P) node[midway,below left] {$\vect{R}$};
    \draw[thick] (P1)--(P2);
    \fill (P1) circle (1.5pt) node[above] {$P_1$};
    \fill (P) circle (1.5pt) node[above] {$P$};
    \fill (P2) circle (1.5pt) node[right] {$P_2$};
    \node at (2.8,2.35) {$m_2$}; \node at (4.9,1.75) {$m_1$};
  \end{tikzpicture}
\end{center}
\end{example}

\subsection{内積（スカラー積）}

\begin{definition}
二つのベクトル $\vect{A},\vect{B}$ のなす角を $\theta$ とするとき、
\begin{equation}
  \vect{A}\cdot\vect{B}=|\vect{A}|\,|\vect{B}|\cos\theta
\end{equation}
を両者の\textbf{内積}という。内積の結果はスカラーである。
\end{definition}

特に $\vect{A}=\vect{B}$ とすれば
\begin{equation}
  \vect{A}\cdot\vect{A}=|\vect{A}|^2,\qquad
  |\vect{A}|=\sqrt{\vect{A}\cdot\vect{A}}
\end{equation}
である。正規直交基底を用いると
\begin{equation}
  \vect{A}\cdot\vect{B}=A_xB_x+A_yB_y+A_zB_z,\qquad
  |\vect{A}|=\sqrt{A_x^2+A_y^2+A_z^2}
\end{equation}
となる。

この成分公式は、基底の直交性から直接導ける。実際、
\begin{align}
 \vect{A}\cdot\vect{B}
 &= (A_x\unit{x}+A_y\unit{y}+A_z\unit{z})
    \cdot(B_x\unit{x}+B_y\unit{y}+B_z\unit{z}) \notag\\
 &=A_xB_x(\unit{x}\cdot\unit{x})
   +A_yB_y(\unit{y}\cdot\unit{y})
   +A_zB_z(\unit{z}\cdot\unit{z}) \notag\\
 &=A_xB_x+A_yB_y+A_zB_z.
\end{align}
異なる基底ベクトルどうしの内積は $0$ なので、交差項はすべて消える。

\begin{figure}[htbp]
  \centering
  \begin{tikzpicture}[>=Stealth]
    \coordinate (O) at (0,0); \coordinate (B) at (4,0); \coordinate (A) at (2.6,2.1);
    \draw[->,very thick] (O)--(B) node[right] {$\vect{B}$};
    \draw[->,very thick] (O)--(A) node[above] {$\vect{A}$};
    \draw[dashed] (A)--(2.6,0);
    \draw[very thick] (O)--(2.6,0) node[midway,below] {$|\vect{A}|\cos\theta$};
    \pic[draw,->,"$\theta$",angle radius=9mm,angle eccentricity=1.25]
      {angle=B--O--A};
  \end{tikzpicture}
  \caption{$\vect{A}$ の $\vect{B}$ 方向への射影}
\end{figure}

内積は交換法則と分配法則を満たす。
\begin{align}
  \vect{A}\cdot\vect{B}&=\vect{B}\cdot\vect{A},\\
  \vect{A}\cdot(\vect{B}+\vect{C})
  &=\vect{A}\cdot\vect{B}+\vect{A}\cdot\vect{C},\\
  a(\vect{A}\cdot\vect{B})&=(a\vect{A})\cdot\vect{B}
  =\vect{A}\cdot(a\vect{B}).
\end{align}

\begin{example}
$P_1=(1,1),P_2=(5,1),P_3=(5,3)$ とする。
$\vect{A}=\overrightarrow{P_1P_3}=(4,2)^{\mathsf T}$、
$\vect{B}=\overrightarrow{P_1P_2}=(4,0)^{\mathsf T}$ なので
\begin{equation}
 \cos\theta=\frac{\vect{A}\cdot\vect{B}}{|\vect{A}|\,|\vect{B}|}
 =\frac{16}{\sqrt{20}\cdot4}=\frac{2}{\sqrt5}.
\end{equation}
\begin{center}
  \begin{tikzpicture}[>=Stealth,scale=.75]
    \draw[->] (0,0)--(6,0) node[right] {$x$};
    \draw[->] (0,0)--(0,4) node[above] {$y$};
    \coordinate (P1) at (1,1); \coordinate (P2) at (5,1);
    \coordinate (P3) at (5,3);
    \draw[->,very thick] (P1)--(P2) node[midway,below] {$\vect{B}$};
    \draw[->,very thick] (P1)--(P3) node[midway,above left] {$\vect{A}$};
    \draw[dashed] (P2)--(P3);
    \fill (P1) circle (2pt) node[left] {$P_1(1,1)$};
    \fill (P2) circle (2pt) node[right] {$P_2(5,1)$};
    \fill (P3) circle (2pt) node[right] {$P_3(5,3)$};
    \pic[draw,->,"$\theta$",angle radius=7mm,angle eccentricity=1.3]
      {angle=P2--P1--P3};
  \end{tikzpicture}
\end{center}
\end{example}

\subsection{外積（ベクトル積）}

\begin{definition}
三次元ベクトル $\vect{A},\vect{B}$ に対して、
\begin{equation}
  \vect{A}\times\vect{B}
  =|\vect{A}|\,|\vect{B}|\sin\theta\,\vect{n}
\end{equation}
を\textbf{外積}という。$\vect{n}$ は $\vect{A},\vect{B}$ の張る面に
垂直な単位ベクトルで、その向きは右ねじの規則で定める。
\end{definition}

$|\vect{A}\times\vect{B}|$ は二つのベクトルが張る平行四辺形の面積に等しい。

\begin{figure}[htbp]
  \centering
  \begin{tikzpicture}[>=Stealth]
    \coordinate (O) at (0,0); \coordinate (A) at (3.4,.5);
    \coordinate (B) at (1.3,1.8); \coordinate (C) at ($(A)+(B)$);
    \fill[blue!8] (O)--(A)--(C)--(B)--cycle;
    \draw[->,very thick] (O)--(A) node[midway,below] {$\vect{A}$};
    \draw[->,very thick] (O)--(B) node[midway,left] {$\vect{B}$};
    \draw (A)--(C)--(B);
    \draw[->,very thick] (O)--(0,3) node[above] {$\vect{A}\times\vect{B}$};
  \end{tikzpicture}
  \caption{外積の向きと平行四辺形の面積}
\end{figure}

外積には次の性質がある。
\begin{align}
  \vect{A}\times\vect{B}&=-\vect{B}\times\vect{A},\\
  \vect{A}\times(\vect{B}+\vect{C})
  &=\vect{A}\times\vect{B}+\vect{A}\times\vect{C},\\
  a(\vect{A}\times\vect{B})&=(a\vect{A})\times\vect{B}
  =\vect{A}\times(a\vect{B}).
\end{align}
正規直交基底では
\begin{equation}
 \unit{x}\times\unit{y}=\unit{z},\quad
 \unit{y}\times\unit{z}=\unit{x},\quad
 \unit{z}\times\unit{x}=\unit{y}
\end{equation}
であり、同じ基底ベクトルどうしの外積は零ベクトルである。したがって
\begin{equation}
 \vect{A}\times\vect{B}=
 \begin{vmatrix}
   \unit{x}&\unit{y}&\unit{z}\\
   A_x&A_y&A_z\\
   B_x&B_y&B_z
 \end{vmatrix}
 =\begin{pmatrix}
 A_yB_z-A_zB_y\\A_zB_x-A_xB_z\\A_xB_y-A_yB_x
 \end{pmatrix}.
\end{equation}

上の行列式を展開すると、各成分がどこから現れるかを確認できる。
\begin{align}
 \vect{A}\times\vect{B}
 ={}&(A_yB_z-A_zB_y)\unit{x}
 +(A_zB_x-A_xB_z)\unit{y} \notag\\
 &+(A_xB_y-A_yB_x)\unit{z}.
\end{align}
外積は順序を入れ替えると符号が反転するため、内積とは異なり交換法則を
満たさない点に注意する。

\begin{example}
$\vect{A}=(1,-1,2)^{\mathsf T}$、$\vect{B}=(3,1,-1)^{\mathsf T}$ のとき
\begin{equation}
 \vect{A}\times\vect{B}
 =\begin{pmatrix}(-1)(-1)-2\cdot1\\2\cdot3-1(-1)\\1\cdot1-(-1)3\end{pmatrix}
 =\begin{pmatrix}-1\\7\\4\end{pmatrix}.
\end{equation}
\end{example}

\subsection{外積による多角形の面積}

$xy$ 平面上の頂点 $P_k=(x_k,y_k,0)^{\mathsf T}$
（$k=1,\ldots,N$）を反時計回りに並べ、$P_{N+1}=P_1$ とする。
原点と隣り合う二頂点が作る符号付き三角形の面積を足すと
\begin{equation}
 S=\frac12\sum_{k=1}^{N}(\vect{P}_k\times\vect{P}_{k+1})_z
  =\frac12\sum_{k=1}^{N}(x_ky_{k+1}-x_{k+1}y_k)
\end{equation}
を得る。頂点が時計回りなら右辺は負になるため、一般には絶対値を取る。
凹多角形でも、外側へ張り出した三角形と切り欠き側の三角形が逆符号に
なるので、同じ公式をそのまま利用できる。

\begin{figure}[htbp]
  \centering
  \begin{tikzpicture}[>=Stealth,scale=.9]
    \draw[->] (-.4,0)--(5.2,0) node[right] {$x$};
    \draw[->] (0,-.4)--(0,4.5) node[above] {$y$};
    \coordinate (P1) at (1,1); \coordinate (P2) at (3.5,.7);
    \coordinate (P3) at (2.7,2.0); \coordinate (P4) at (4.4,2.9);
    \coordinate (P5) at (2.8,3.8); \coordinate (P6) at (.8,3);
    \fill[blue!8] (P1)--(P2)--(P3)--(P4)--(P5)--(P6)--cycle;
    \draw[very thick] (P1)--(P2)--(P3)--(P4)--(P5)--(P6)--cycle;
    \foreach \p/\name in {P1/P_1,P2/P_2,P3/P_3,P4/P_4,P5/P_5,P6/P_6}
      {\fill (\p) circle (1.6pt) node[above right] {$\name$};\draw[gray] (0,0)--(\p);}
  \end{tikzpicture}
  \caption{三角形への分割による多角形の符号付き面積}
\end{figure}

\begin{example}
$P_1=(1,2),P_2=(5,2),P_3=(5,4)$ の三角形では
\begin{equation}
 S=\frac12\{(1\cdot2-5\cdot2)+(5\cdot4-5\cdot2)
 +(5\cdot2-1\cdot4)\}=4.
\end{equation}
\end{example}
